% !TEX root = my-thesis.tex

% **************************************************
% Information and Commands for Reuse
% **************************************************
\newcommand{\thesisTitle}{Structuring LLM Arguments:
A Logic of Argumentation Based Approach}
\newcommand{\thesisName}{Aaron Habyarimana}
\newcommand{\thesisBirthplace}{Mainz}
% Options: Bachelorarbeit, Masterarbeit, Bachelor thesis, Master thesis, Dissertation.
% Dissertation has different titlepage than others
\newcommand{\thesisSubject}{Bachelorarbeit}
\newcommand{\thesisKeywords}{Abstract Argumentation Frameworks, Argument Mining, Dung Semantics, Large Language Models, Formal Reasoning}  % (optional) Keywords are added to PDF meta data

\newcommand{\thesisDate}{July 15, 2026} % submission date (fixed, do not use \today)
\ifcsname thesisVersion\endcsname
% version specifier from cmd option. In the gitlab CI/CD job the commit short SHA will be inserted.
\else
% update to fixed version string if no command option is specified
\newcommand{\thesisVersion}{CleanThesis 0.4.1 - JGU patches}
\fi

\newcommand{\thesisFirstReviewer}{Prof. Dr. Stefan Kramer}
\newcommand{\thesisFirstReviewerUniversity}{\protect{Johannes Gutenberg University Mainz}}
\newcommand{\thesisFirstReviewerDepartment}{Institute of Computer Science - Data Mining}

\newcommand{\thesisSecondReviewer}{Dr. Maike Paetzel-Prüsmann}
\newcommand{\thesisSecondReviewerUniversity}{\protect{Johannes Gutenberg University Mainz}}
\newcommand{\thesisSecondReviewerDepartment}{Institute of Computer Science}

\newcommand{\thesisFirstSupervisor}{Derian Boer}
\newcommand{\thesisSecondSupervisor}{Lennart Baur}

\newcommand{\thesisUniversity}{\protect{Johannes Gutenberg University Mainz}}
\newcommand{\thesisUniversityDepartment}{FB08}
\newcommand{\thesisUniversityDepartmentLong}{Faculty of Physics, Mathematics, and Computer Science}
\newcommand{\thesisUniversityInstitute}{Institute of Computer Science}
\newcommand{\thesisUniversityGroup}{Data Mining}
\newcommand{\thesisUniversityCity}{Mainz}
\newcommand{\thesisUniversityStreetAddress}{Staudingerweg 9}
\newcommand{\thesisUniversityPostalCode}{55128}

\newcommand{\thesisAIToolsAllowed}{true}  % Change here [true, false] for declaration.
\newcommand{\thesisTitlepageElement}{show}  % show or hide titlepage graphical element
\newcommand{\thesisTitlepageHeight}{6cm}  % height of line in cm
\newcommand{\thesisTitlepageStart}{8cm}  % shifted start from center


% **************************************************
% Debug LaTeX Information
% **************************************************
%\listfiles


% **************************************************
% Load and Configure Packages
% **************************************************
\usepackage[ngerman, english]{babel} % babel system, adjust the language of the content (the last mentioned language is used as default for all text)
\PassOptionsToPackage{% setup clean thesis style
    figuresep=colon,%
    hangfigurecaption=false,%
    hangsection=true,%
    hangsubsection=true,%
    sansserif=false,%
    configurelistings=true,%
    colorize=full,%
    colortheme=jgured,%
    configurebiblatex=true,%
    bibsys=biber,%
    bibfile=bib-refs,%
    bibstyle=numeric,%
    bibsorting=nty,%
}{cleanthesis}
\usepackage{cleanthesis}

\hypersetup{% setup the hyperref-package options
    pdftitle={\thesisTitle},    %   - title (PDF meta)
    pdfsubject={\thesisSubject},%   - subject (PDF meta)
    pdfauthor={\thesisName},    %   - author (PDF meta)
    pdfkeywords={\thesisKeywords}, % list of keywords
    plainpages=false,           %   -
    colorlinks=false,           %   - colorize links?
    pdfborder={0 0 0},          %   -
    breaklinks=true,            %   - allow line break inside links
    bookmarksnumbered=true,     %
    bookmarksopen=true          %
}

% **************************************************
% Other Packages
% **************************************************
\usepackage{scrhack}
\usepackage{iflang}  % for language switchcase in titlepage
\usepackage{tikz}
\usetikzlibrary{arrows.meta, positioning, shapes.geometric, fit, backgrounds}
\usepackage{amsmath}
\usepackage{amssymb}
\usepackage{amsthm}
\usepackage{booktabs}
\usepackage{algorithm}
\usepackage{algpseudocode}

% theorem-like environments used in the foundations chapter
\theoremstyle{definition}
\newtheorem{definition}{Definition}[chapter]
\newtheorem{proposition}[definition]{Proposition}
